\pdfvariable suppressoptionalinfo 512\relax
\pdfvariable trailerid{}
\documentclass[10pt]{article}
\usepackage[a4paper,margin=0.68in]{geometry}
\usepackage{fontspec}
\usepackage{amsmath,amssymb,booktabs}
\newfontfamily\cjkfont{Droid Sans Fallback}
\linespread{0.96}
\title{\vspace{-1.5em}Renormalized Branch Amplitudes for the Square-Billiard Abel Trace}
\author{Route-A structural certificate C162}
\date{}
\begin{document}
\maketitle
\begin{abstract}
For the Dirichlet half-wave trace $W_D$ of the unit square, we prove that every
nonzero source lattice shell admits a canonical full-trace normalization.  At
$t_N=2\sqrt N$,
$\epsilon^{3/2}W_D(\epsilon-it_N)$ tends to
$e^{i\pi/4}r_2^{\rm src}(N)/(8\pi N^{1/4})$; negative time gives its complex
conjugate.  This advances the earlier branch-location theorem by extracting a
finite coefficient from the complete trace, including times where the
boundary subtraction has a coincident simple pole.  The proof fixes the phase
with the principal branch, splits off finitely many shells, and uniformly
dominates the infinite remainder by the summable two-dimensional
$|m|^{-3}$ tail.  Exact shells through $N=800$ and high-precision convergence
rows are regression sentinels, not the proof.  The coefficient is aggregate
clean-family multiplicity, not an isolated stability amplitude or target
divisor law, and no arithmetic Euler or Hilbert--P\'olya claim is made.
\end{abstract}
\noindent\textbf{Keywords:} square billiard; Abel half-wave trace; clean
family; branch singularity; boundary normalization; Poisson summation.

\medskip
\noindent{\cjkfont\parbox{\linewidth}{\raggedright\emergencystretch=2em
中\hspace{0pt}文\hspace{0pt}摘\hspace{0pt}要\hspace{0pt}：本\hspace{0pt}文\hspace{0pt}为\hspace{0pt}单\hspace{0pt}位\hspace{0pt}正\hspace{0pt}方\hspace{0pt}形\hspace{0pt}狄\hspace{0pt}利\hspace{0pt}克\hspace{0pt}雷\hspace{0pt}半\hspace{0pt}波\hspace{0pt}迹\hspace{0pt}的\hspace{0pt}每\hspace{0pt}个\hspace{0pt}非\hspace{0pt}零\hspace{0pt}源\hspace{0pt}格\hspace{0pt}点\hspace{0pt}壳\hspace{0pt}层\hspace{0pt}建\hspace{0pt}立\hspace{0pt}规\hspace{0pt}范\hspace{0pt}的\hspace{0pt}
全\hspace{0pt}迹\hspace{0pt}边\hspace{0pt}界\hspace{0pt}归\hspace{0pt}一\hspace{0pt}化\hspace{0pt}。\hspace{0pt}\allowbreak 在\hspace{0pt} $t_N=2\sqrt N$ 处\hspace{0pt}，
$\epsilon^{3/2}W_D(\epsilon-it_N)$ 的\hspace{0pt}极\hspace{0pt}限\hspace{0pt}为\hspace{0pt}
$e^{i\pi/4}r_2^{\rm src}(N)/(8\pi N^{1/4})$，\allowbreak
负\hspace{0pt}时\hspace{0pt}间\hspace{0pt}极\hspace{0pt}限\hspace{0pt}为\hspace{0pt}其\hspace{0pt}复\hspace{0pt}共\hspace{0pt}轭\hspace{0pt}。\hspace{0pt}\allowbreak 该\hspace{0pt}结\hspace{0pt}论\hspace{0pt}不\hspace{0pt}仅\hspace{0pt}定\hspace{0pt}位\hspace{0pt}分\hspace{0pt}支\hspace{0pt}点\hspace{0pt}，而\hspace{0pt}且\hspace{0pt}从\hspace{0pt}完\hspace{0pt}整\hspace{0pt}迹\hspace{0pt}中\hspace{0pt}
提\hspace{0pt}取\hspace{0pt}有\hspace{0pt}限\hspace{0pt}系\hspace{0pt}数\hspace{0pt}；\allowbreak 即\hspace{0pt}便\hspace{0pt}边\hspace{0pt}界\hspace{0pt}扣\hspace{0pt}除\hspace{0pt}项\hspace{0pt}在\hspace{0pt}同\hspace{0pt}一\hspace{0pt}时\hspace{0pt}刻\hspace{0pt}有\hspace{0pt}简\hspace{0pt}单\hspace{0pt}极\hspace{0pt}点\hspace{0pt}，
归\hspace{0pt}一\hspace{0pt}化\hspace{0pt}后\hspace{0pt}该\hspace{0pt}极\hspace{0pt}点\hspace{0pt}仍\hspace{0pt}消\hspace{0pt}失\hspace{0pt}。\hspace{0pt}\allowbreak 证\hspace{0pt}明\hspace{0pt}以\hspace{0pt}主\hspace{0pt}值\hspace{0pt}分\hspace{0pt}支\hspace{0pt}固\hspace{0pt}定\hspace{0pt}相\hspace{0pt}位\hspace{0pt}，再\hspace{0pt}分\hspace{0pt}离\hspace{0pt}有\hspace{0pt}限\hspace{0pt}壳\hspace{0pt}层\hspace{0pt}，
并\hspace{0pt}以\hspace{0pt}二\hspace{0pt}维\hspace{0pt} $|m|^{-3}$ 可\hspace{0pt}求\hspace{0pt}和\hspace{0pt}尾\hspace{0pt}项\hspace{0pt}统\hspace{0pt}一\hspace{0pt}控\hspace{0pt}制\hspace{0pt}无\hspace{0pt}限\hspace{0pt}余\hspace{0pt}项\hspace{0pt}。\hspace{0pt}\allowbreak
直\hspace{0pt}到\hspace{0pt} $N=800$ 的\hspace{0pt}精\hspace{0pt}确\hspace{0pt}壳\hspace{0pt}层\hspace{0pt}与\hspace{0pt}高\hspace{0pt}精\hspace{0pt}度\hspace{0pt}数\hspace{0pt}值\hspace{0pt}仅\hspace{0pt}作\hspace{0pt}回\hspace{0pt}归\hspace{0pt}哨\hspace{0pt}兵\hspace{0pt}。\hspace{0pt}\allowbreak
所\hspace{0pt}得\hspace{0pt}系\hspace{0pt}数\hspace{0pt}表\hspace{0pt}示\hspace{0pt}光\hspace{0pt}滑\hspace{0pt}轨\hspace{0pt}道\hspace{0pt}族\hspace{0pt}总\hspace{0pt}重\hspace{0pt}数\hspace{0pt}，不\hspace{0pt}是\hspace{0pt}孤\hspace{0pt}立\hspace{0pt}轨\hspace{0pt}道\hspace{0pt}稳\hspace{0pt}定\hspace{0pt}性\hspace{0pt}振\hspace{0pt}幅\hspace{0pt}，
也\hspace{0pt}不\hspace{0pt}推\hspace{0pt}出\hspace{0pt}目\hspace{0pt}标\hspace{0pt}除\hspace{0pt}数\hspace{0pt}律\hspace{0pt}、\hspace{0pt}欧\hspace{0pt}拉\hspace{0pt}数\hspace{0pt}据\hspace{0pt}或\hspace{0pt}希\hspace{0pt}尔\hspace{0pt}伯\hspace{0pt}特\hspace{0pt}波\hspace{0pt}利\hspace{0pt}亚\hspace{0pt}算\hspace{0pt}子\hspace{0pt}。\hspace{0pt}\par\smallskip
关\hspace{0pt}键\hspace{0pt}词\hspace{0pt}：正\hspace{0pt}方\hspace{0pt}形\hspace{0pt}台\hspace{0pt}球\hspace{0pt}；\allowbreak 阿\hspace{0pt}贝\hspace{0pt}尔\hspace{0pt}半\hspace{0pt}波\hspace{0pt}迹\hspace{0pt}；\allowbreak 光\hspace{0pt}滑\hspace{0pt}轨\hspace{0pt}道\hspace{0pt}族\hspace{0pt}；
\allowbreak 分\hspace{0pt}支\hspace{0pt}奇\hspace{0pt}点\hspace{0pt}；\allowbreak 边\hspace{0pt}界\hspace{0pt}归\hspace{0pt}一\hspace{0pt}化\hspace{0pt}；\allowbreak 泊\hspace{0pt}松\hspace{0pt}求\hspace{0pt}和\hspace{0pt}。\hspace{0pt}
}}

\section{Exact source trace}
For $\Re s>0$, the unit-square Dirichlet trace is
\[
W_D(s)=\operatorname{Tr}e^{-s\sqrt{\Delta_D}}
=\sum_{j,k\ge1}e^{-\pi s\sqrt{j^2+k^2}}.
\]
C157 established the exact principal-branch Poisson representation
\[
W_D(s)=\frac{s}{2\pi}\sum_{m\in\mathbb Z^2}
(s^2+4|m|^2)^{-3/2}-\frac14-\frac1{e^{\pi s}-1}.              \tag{1}
\]
The dual series is locally normally convergent in the right half-plane.

\section{Canonical full-trace limit}
Let $N\ge1$, $t_N=2\sqrt N$, and
$r_2^{\rm src}(N)=\#\{m\in\mathbb Z^2:|m|^2=N\}$.  Then
\[
\boxed{\lim_{\epsilon\downarrow0}\epsilon^{3/2}
W_D(\epsilon-it_N)=
\frac{e^{i\pi/4}r_2^{\rm src}(N)}{8\pi N^{1/4}}.}            \tag{2}
\]
At $-t_N$ the limit is the complex conjugate.  Indeed the spectral
coefficients are real, so
$W_D(\overline{s})=\overline{W_D(s)}$ throughout the right half-plane; the
principal dual representation in (1) has the same conjugation symmetry.

For a matching vector, with $s=\epsilon-it_N$,
$s^2+4N=\epsilon(\epsilon-2it_N)$.  Its normalized contribution tends to
\[
\frac{-it_N}{2\pi}(-2it_N)^{-3/2}
=\frac{e^{i\pi/4}}{8\pi N^{1/4}},                            \tag{3}
\]
where the phase is fixed by the principal power.  Multiplication by the number
of matching vectors gives the right side of (2).

The infinite remainder requires a uniform argument.  Fix $t_N$ and
$0<\epsilon\le1$.  Choose $R$ so that
$4|m|^2\ge2(t_N^2+1)$ for $|m|\ge R$.  Then
\[
|s^2+4|m|^2|\ge4|m|^2-|s|^2\ge2|m|^2,
\]
so the tail is dominated uniformly by a constant times
$\sum_{m\ne0}|m|^{-3}<\infty$.  The finitely many nonmatching shells stay
bounded.  Dominated convergence and $\epsilon^{3/2}$ kill all these terms.
If $t_N$ is even, the subtraction term in (1) can have a simple pole, but its
normalized size is only $O(\epsilon^{1/2})$.  More precisely, this coincidence
means $N=k^2$ and $t_N=2k$, whence
\[
-\frac1{e^{\pi(\epsilon-i2k)}-1}
=-\frac1{\pi\epsilon}+\frac12+O(\epsilon).
\]
After multiplication by $\epsilon^{3/2}$ the pole is
$-\epsilon^{1/2}/\pi+O(\epsilon^{3/2})$ and vanishes.  The Weyl zero-mode is
bounded at $t_N\ne0$, and the constant term is also killed.  This proves a
limit for the full trace, not just one summand.

\section{Certificate and boundary}
Exact reconstruction through $N=800$ finds 270 occupied shells and 2520
nonzero lattice vectors; 28 occupied square norms coincide with simple-pole
times.  The shell $N=65$ retains the first four ordered positive primitive
directions.  Local high-precision rows at $N=1,2,5,13,65$ check convergence
but are not the proof.
For each shell, the exact ledger separately records the four axis vectors when
$N$ is a square and all nonaxis sign lifts; gcd reduction of every nonaxis
vector retains its primitive direction and repetition.  Their total is exactly
$r_2^{\rm src}(N)$ in (2), so no multiplicity is inferred from floating data.

The strict tuple is $(\texttt{A1\_WEAK},\texttt{A2\_FAIL},
\texttt{A3\_FAIL},\texttt{A4\_NATURAL\_QUANTIZATION})$:
$\sqrt{\Delta_D}$ is a natural self-adjoint source operator, but the
coefficient aggregates positive-dimensional clean families.  We claim no
isolated stability determinant, target trace/divisor/counting law, arithmetic
local/Euler factor, root number, automorphy, Hilbert--P\'olya construction, or
Route-B authorization.  Scope: \texttt{NO\_BAD\_EULER\_OR\_ROOT\_NUMBER}.

\section*{Declarations}
\textbf{Data availability.} All claim-bearing data and code are included; no
external dataset is used.  \textbf{Ethics.} No human participants, animals, or
sensitive personal data are involved.  \textbf{Author contributions/CRediT.}
Anonymous technical certificate; Conceptualization, Formal analysis, Software,
Validation, and Writing are recorded at package level.  \textbf{Competing
interests.} None known.  \textbf{Funding.} No external funding is reported.
\textbf{AI-use disclosure.} An AI coding assistant supported derivation
planning, drafting, and code review.  Packaged independent checks validate
quantitative claims; the assistant was not treated as an external peer reviewer.
\end{document}
